\chapter*{Before you start}
\addcontentsline{toc}{chapter}{Before you start}
\markboth{Before you start}{Before you start}

This guide was written from scratch, with original explanations, to help you learn cross-section econometrics, panel data and time series. It is \emph{not} a translation of any of the books. The books are the source for deeper reading. This material is the didactic path that leads you to them.

\section*{How the guide is organized}
Each chapter follows the same order: the idea in plain language, the essential mathematics, a numerical experiment that confirms the theory, the most common pitfalls, and exercises with solutions. The colored boxes have fixed roles:

\begin{intuicao}[Intuition]
Explains the \emph{why} before the formula.
\end{intuicao}
\begin{atencao}[Pitfall]
Frequent mistakes, including some that appear in books and software manuals.
\end{atencao}
\begin{simulacao}[Checked by simulation]
A result I ran on the computer (Python) to make sure the statement is correct. The numbers come from a fixed seed, so they vary slightly if you repeat the experiment.
\end{simulacao}

\section*{Map of the books: where to go deeper}
\begin{center}
\small
\begin{tabular}{p{4.3cm}p{9.6cm}}
\toprule
\textbf{Topic} & \textbf{Where to read more} \\
\midrule
Background probability and linear algebra & Stachurski (\emph{Econometric Theory}, lecture notes), parts I and II; Goldberger \\
Regression, OLS, IV & Wooldridge (ch.\ 4 and 5 of the 1st edition); Hanck et al.\ (ch.\ 4 to 9 and 12) \\
Panel data & Wooldridge (ch.\ 10 and 11); Hanck et al.\ (ch.\ 10) \\
ARMA, ARIMA, Box--Jenkins & Cryer and Chan (ch.\ 4 to 9); Mills (ch.\ 3, 4 and 7); Hamilton (ch.\ 3 and 4) \\
Unit roots and cointegration & Hamilton (ch.\ 17 to 19); Mills (ch.\ 14); Hanck et al.\ (ch.\ 14 and 16) \\
VAR & Lütkepohl (\emph{New Introduction to Multiple Time Series Analysis}); Hamilton (ch.\ 11); Mills (ch.\ 13) \\
Volatility (ARCH/GARCH) & Tsay (\emph{Analysis of Financial Time Series}); Mills (ch.\ 10); Hamilton (ch.\ 21) \\
Machine learning & Hastie, Tibshirani and Friedman (\emph{Elements of Statistical Learning}); Faloutsos et al.\ tutorial on large-scale forecasting; Faouzi and Janati (\texttt{pyts} package) \\
\bottomrule
\end{tabular}
\end{center}

\section*{Two remarks about the books}
\textbf{Wooldridge edition.} The file you sent is the 1st edition of \emph{Econometric Analysis of Cross Section and Panel Data} (2002, 2004 printing). A 2nd edition appeared in 2010 and is considerably longer. The concepts I use here (OLS, IV, fixed and random effects) are in both, but chapter numbering and some advanced topics change. For current applied work, the 2nd edition is worth having.

\textbf{About ``correcting the books' errors''.} I did not audit thousands of pages looking for errata, and I will not pretend I did. What I did was (i) check by simulation every formula and number that appears in this guide and (ii) highlight, in the \emph{Pitfall} boxes, the points where books disagree with each other or where a mistake is common. The most important one is the sign of the moving-average coefficient, which changes from book to book (chapter~\ref{cap:est}). I also did not read the books cover to cover: I used the tables of contents to build the map above and checked directly in Cryer--Chan and in Mills the sign convention of the MA. If you find something that does not match your book, trust the simulation and the derivation, not my text.

\section*{Notation}
Vectors and matrices are in bold when useful. $\E(\cdot)$ is expectation, $\Var$ variance, $\Cov$ covariance, $\plim$ probability limit. $\hat\beta$ is an estimator and $\beta$ the true parameter. In time series, $B$ is the lag operator: $B y_t = y_{t-1}$.
